% Ad Familiares 16.26 --- headnote
% ad-familiares-16-26, c. July 44 BC, place uncertain

\textit{Quintus Cicero --- Marcus's younger brother --- to Tiro,
Marcus's freedman and secretary. One of only a handful of letters in
the Ciceronian corpus written by Quintus rather than by Marcus; the
salutation \textit{QVINTVS TIRONI SVO} makes the sender explicit.
The Perseus dateline places it ``in an uncertain place and month of
710'' (44 BC), the year of the Ides of March and its aftermath.}

\textit{The note is brief and warm. Quintus has received a second
packet of mail without a letter from Tiro and is scolding him for
it --- mock-judicially, since with himself as advocate Tiro cannot
hope to be acquitted and would need to call in Marcus as defense
counsel. The closing image, drawn from family memory, has the
quality of an unguarded household anecdote: their mother used to
seal even empty wine-flasks, so that no one could claim a flask had
been empty all along when it had in fact been quietly drained. Tiro
should write even when he has nothing to write about, so as not to
``steal the time.'' The plea ends ``love us'' --- a brother's
affection extended through Marcus's household to his freedman.}
